\section{Numerical method and mathematical observables}

\subsection{Overview of the calculation map}

For one request $(N,\rho,\{a_r\},\text{gates},\text{tolerances})$, the production
calculation performs the composition
\begin{equation}
 K_\rho(q_j)
 \longrightarrow (\lambda_n(q_j),V_j)
 \longrightarrow F_j
 \longrightarrow H_j
 \longrightarrow A_{a_r,j}
 \longrightarrow
 \begin{cases}
  \text{restricted-alignment observables},\\
  \text{cyclic Fourier powers},\\
  \text{ordinary truncation consequences}.
 \end{cases}
 \label{eq:calculation-map}
\end{equation}
The parent is assembled explicitly, diagonalized by a Hermitian eigensolver, and then
qualified before any gauge-amplitude point is evaluated.  Parent preparation is reused
within an amplitude study so that changing $a$ cannot silently change the parent,
transported frame, or reference Fourier field.

\subsection{Closed-loop polar frame transport}

Let $V_j\in\Complex^{4\times2}$ be a raw orthonormal spectral frame.  Neighboring frames
are related through overlap matrices
\begin{equation}
 M_j=F_{j-1}^\dagger V_j\in\Complex^{2\times2}.
 \label{eq:neighbor-overlap}
\end{equation}
For a singular-value decomposition $M_j=L_j\Sigma_jR_j^\dagger$, the unitary polar
factor is $Q_j=L_jR_j^\dagger$.  A right action based on $Q_j$ chooses the member of the
same retained subspace that is closest to the preceding transported frame.  The minimum
singular value over all links diagnoses whether this polar step is numerically resolved.
The installed request uses a rank-relative tolerance to decide numerical rank and a
separate validation-relative tolerance for backend residual checks.

Transport around the periodic loop generally leaves an endpoint mismatch
$W\in U(2)$.  With a Schur or spectral representation
\begin{equation}
 W=Z\diag(e^{i\varphi_1},e^{i\varphi_2})Z^\dagger,
 \qquad -\pi<\varphi_n\le\pi,
 \label{eq:closure-holonomy}
\end{equation}
principal-branch phases are distributed along the loop.  At mesh fraction $j/N$, the
correction has the form
\begin{equation}
 C_j=Z\diag\!\left(e^{-ij\varphi_1/N},e^{-ij\varphi_2/N}\right)Z^\dagger.
 \label{eq:closure-distribution}
\end{equation}
The resulting frame is periodic under the declared identity endpoint sewing.  This is a
frame construction, not a change of spectral subspace.

Qualification retains, rather than hides, the following diagnostics:
\begin{align}
 \sigma_{\min}&=\min_{j,n}\sigma_n(M_j),\\
 d_{\mathrm{unitary}}&=\max_j\norm{Q_j^\dagger Q_j-I_2}_{\max},\\
 d_{\mathrm{closure}}&=\norm{\text{computed closure}-\text{declared closure}}_{\max},\\
 d_F&=\max_j\norm{F_j^\dagger F_j-I_2}_F,\\
 d_H&=\max_j\norm{H_j-H_j^\dagger}_{\max}.
 \label{eq:transport-diagnostics}
\end{align}
A backend exception, numerical rank failure, policy threshold failure, external-gap
failure, and successful qualification are five different outcomes.  They are not
encoded as a single boolean.

\subsection{Restricted global alignment}

The comparison does not admit the local gauge $\mathcal R(\theta_a(q))$ as a recovery
map.  Instead it admits one global real rotation
\begin{equation}
 G(\phi)=\mathcal R(\phi),\qquad \phi\in[0,\pi),
 \label{eq:global-rotation}
\end{equation}
where the interval has length $\pi$ because conjugation of a traceless rank-two matrix
has a doubled-angle periodicity.

For discrete traceless fields $X_j$ and $Y_{a,j}$, define
\begin{equation}
 c_{N,a}(\phi)
 =\frac{1}{2N}\sum_{j=0}^{N-1}
 \Re\tr\!\left[X_j^\dagger G(\phi)^\dagger
                  Y_{a,j}G(\phi)\right].
 \label{eq:operator-correlation}
\end{equation}
Rank-two real rotations imply an exact second-harmonic representation
\begin{equation}
 c_{N,a}(\phi)=c_0+c_c\cos(2\phi)+c_s\sin(2\phi).
 \label{eq:correlation-harmonic}
\end{equation}
The coefficients are formed directly from the matrix fields, and a continuous harmonic
maximizer returns
\begin{equation}
 c_{N,a}^{\max}=\max_\phi c_{N,a}(\phi),
 \qquad
 \phi_{N,a}^{\star}\in\operatorname*{arg\,max}_\phi c_{N,a}(\phi).
 \label{eq:harmonic-maximum}
\end{equation}
The maximizing argument can be absent when the harmonic is degenerate under the declared
absolute and relative degeneracy tolerances.  An absent argument is a scientific
outcome, not the number zero.

The normalized discrete maximum correlation is
\begin{equation}
 \mu_{N,a}=\frac{c_{N,a}^{\max}}{\kappa_N},
 \qquad -1\le\mu_{N,a}\le1,
 \label{eq:normalized-correlation}
\end{equation}
up to bounded binary64 roundoff.  Values just outside the interval by no more than
$32\epsilon_{64}$ are clamped; a larger excursion is a floating-point failure.  The
identity amplitude $a=0$ is installed as the exact invariant $\mu_{N,0}=1$ in the
production route.

\subsection{Two non-equivalent restricted thresholds}

The first comparison allows a nonnegative scalar $s$ together with the global rotation.
Its normalized mean squared loss is
\begin{equation}
 \mathcal L_N(s,\phi;a)
 =\frac{1}{2N}\sum_j
  \norm{Y_{a,j}-sG(\phi)X_jG(\phi)^\dagger}_F^2.
 \label{eq:scaled-loss}
\end{equation}
Because $X$ and $Y_a$ have the same total traceless power $\kappa_N$, minimizing the
quadratic in $s\ge0$ yields
\begin{equation}
 s_{N,a}^{\star}=\max(0,\mu_{N,a}),
 \qquad
 \mathcal L_{N,\min}(a)
 =\kappa_N\left[1-\max(0,\mu_{N,a})^2\right].
 \label{eq:scaled-minimum-loss}
\end{equation}
The retained operator threshold is therefore
\begin{equation}
 \tau_{\mathrm{op},N}(a)
 =\sqrt{\mathcal L_{N,\min}(a)}.
 \label{eq:operator-threshold}
\end{equation}

The second comparison fixes the traceless scale, as required by exact spectrum
preservation, and minimizes only over the global rotation.  Its threshold is
\begin{equation}
 \tau_{\mathrm{joint},N}(a)
 =\sqrt{2\kappa_N[1-\mu_{N,a}]}.
 \label{eq:joint-threshold}
\end{equation}
Equations~\eqref{eq:operator-threshold} and \eqref{eq:joint-threshold} answer different
questions.  The first asks for the best globally rotated and rescaled traceless operator;
the second asks for the best globally rotated operator at the exact spectral scale.  The
report never replaces one by the other.

The production route explicitly reconstructs the matrix loss at
$(s_{N,a}^{\star},\phi_{N,a}^{\star})$ and retains
\begin{equation}
 d_{\mathrm{recon}}(a)
 =\abs{\mathcal L_N(s_{N,a}^{\star},\phi_{N,a}^{\star};a)
       -\mathcal L_{N,\min}(a)}
 \label{eq:loss-reconstruction-defect}
\end{equation}
as a residual channel.

\subsection{Flat-parent analytic formulas}

For the flat traceless control $D$, the gauge rotation changes its orientation through
twice the angle $\theta_a(q)$.  Averaging the aligned correlation over the continuous
Brillouin zone gives
\begin{equation}
 \frac{1}{1}\int_{-1/2}^{1/2}\cos[2a\sin(2\pi q)]\,dq
 =J_0(2a),
 \label{eq:bessel-correlation}
\end{equation}
where $J_0$ is the Bessel function of the first kind.  Consequently,
\begin{align}
 \mu_{\infty,a}^{\mathrm{flat}}&=J_0(2a),\\
 \tau_{\mathrm{op},\infty}^{\mathrm{flat}}(a)
 &=\sqrt{\kappa\,[1-J_0(2a)^2]},\\
 \tau_{\mathrm{joint},\infty}^{\mathrm{flat}}(a)
 &=\sqrt{2\kappa\,[1-J_0(2a)]}.
 \label{eq:flat-thresholds}
\end{align}
These formulas are analytic comparators only for the flat control.  The dispersive
embedded parent does not inherit them.

\subsection{Finite cyclic Fourier representation}

For a sampled matrix field $B_j=B(q_j)$ on Eq.~\eqref{eq:centered-mesh}, the declared
discrete hopping blocks are
\begin{equation}
 B_R=\frac{1}{N}\sum_{j=0}^{N-1}e^{-2\pi iq_jR}B_j,
 \qquad
 B_j=\sum_{R\in\mathcal R_N}e^{2\pi iq_jR}B_R,
 \label{eq:discrete-fourier-pair}
\end{equation}
with a centered cyclic representative inventory $\mathcal R_N$.  The normalized block
power is
\begin{equation}
 p_R[B]=\frac{1}{2}\norm{B_R}_F^2,
 \label{eq:block-power}
\end{equation}
and the non-onsite weight is
\begin{equation}
 W_{\ne0}[B]=\sum_{R\in\mathcal R_N\setminus\{0\}}p_R[B].
 \label{eq:nonzero-weight}
\end{equation}
The reference and comparison weights are retained separately.  Their difference is not
called an error unless a particular comparison explicitly defines one.

The Parseval residual is computed from the same normalization:
\begin{equation}
 d_{\mathrm{Parseval}}[B]
 =\abs{\frac{1}{2N}\sum_j\norm{B_j}_F^2
       -\sum_{R\in\mathcal R_N}p_R[B]}.
 \label{eq:parseval-defect}
\end{equation}
Small residuals establish consistency of the finite transform and normalization.  They
do not establish continuum convergence.

For even $N$, $R=-N/2$ is a centered Nyquist representative that is self-conjugate only
modulo $N$.  Treating it as an ordinary integer displacement would impose the wrong
Hermitian-pairing semantics.  The truncation channel therefore removes this cyclic
Nyquist block, retains its power separately, and converts the remaining inventory to an
ordinary conjugate-closed model satisfying
\begin{equation}
 B_{-R}=B_R^\dagger.
 \label{eq:ordinary-hermiticity}
\end{equation}

\subsection{Parent-hopping family}

The parent-family inventory is
\begin{equation}
 \rho\in\{0,0.25,0.5,0.75,1\}.
 \label{eq:rho-inventory}
\end{equation}
For each $\rho$, the complete procedure---parent diagonalization, gap gate, frame
transport, retained projection, gauge-amplitude sweep, restricted alignment, and Fourier
analysis---is repeated.  It would be incorrect to reuse the $\rho=1$ frame or projected
Hamiltonian at another family member.  The family therefore measures sensitivity to the
ambient parent-hopping homotopy, not merely a scalar interpolation of final observables.

\subsection{Nested momentum refinement}

The production momentum ladder is
\begin{equation}
 N\in\{32,64,128,256,512\},
 \label{eq:momentum-ladder}
\end{equation}
with exact doubling.  The fine mesh contains every coarse point.  At shared points, the
study compares parent eigenvalues and spectral projectors rather than raw eigenvectors,
because raw eigenvectors are gauge dependent.  The shared-point projector defect is
schematically
\begin{equation}
 d_P(N,2N)=\max_{q\in\mathcal Q_N}
 \norm{P_N(q)-P_{2N}(q)}_F.
 \label{eq:projector-refinement}
\end{equation}

For each adjacent pair $(N,2N)$ the retained step contains separate maxima for parent
qualification diagnostics, shared parent eigenvalues, shared projectors, correlation,
operator thresholds, hopping weights, individual hopping powers, maximizing angles,
band defects, loss-reconstruction defects, and Parseval residuals.  Angle changes are
compared only when the maximizing-angle dispositions are compatible; degeneracy status
is retained separately.

Let $d_{r}(N,2N)$ denote any declared change metric and $\eta_r$ its threshold.  A step
passes only if every change metric and every fine-mesh residual satisfies its own bound:
\begin{equation}
 \operatorname{pass}(N,2N)
 =\bigwedge_r[d_r(N,2N)\le\eta_r].
 \label{eq:momentum-pass}
\end{equation}
Convergence requires the declared number of final adjacent pairs to pass.  Earlier
passing pairs cannot compensate for a failing final pair, and one small observable
cannot compensate for another large one.

\subsection{Adaptive amplitude refinement}

The amplitude domain is $[0,\pi/4]$.  Starting from interpolation nodes
$\{a_i\}$, every interval $[a_i,a_{i+1}]$ is first tested at its midpoint
$m_i=(a_i+a_{i+1})/2$.  For a scalar observable $f$, the midpoint linear-interpolation
defect is
\begin{equation}
 d_f^{(1/2)}(i)
 =\abs{f(m_i)-\frac{f(a_i)+f(a_{i+1})}{2}}.
 \label{eq:midpoint-defect}
\end{equation}
The same construction is applied separately to correlation, both thresholds,
comparison non-onsite weight, and every comparison hopping power.  For the hopping-power
channel, the maximum is taken over the full displacement inventory after verifying that
the inventories agree.

Only failing intervals are refined.  If all midpoint checks pass, the node inventory is
frozen and validation moves to disjoint quarter points
\begin{equation}
 v_{i,1}=a_i+\frac14(a_{i+1}-a_i),
 \qquad
 v_{i,3}=a_i+\frac34(a_{i+1}-a_i).
 \label{eq:quarter-points}
\end{equation}
For $\xi\in\{1/4,3/4\}$ the validation defect is
\begin{equation}
 d_f^{(\xi)}(i)
 =\abs{f(a_i+\xi\Delta a_i)
       -[(1-\xi)f(a_i)+\xi f(a_{i+1})]}.
 \label{eq:quarter-defect}
\end{equation}
Convergence requires both a passing midpoint level and a passing disjoint quarter-point
validation on that same frozen node inventory.  A midpoint pass alone is insufficient.
The closed alternatives are qualification stop, converged, refinement-level exhaustion,
and evaluated-amplitude-cap exhaustion.  No alternative is silently relabeled as
convergence.

\subsection{Ordinary-lattice truncation consequences}

After Nyquist removal, let the ordinary hopping model be
\begin{equation}
 B^{\mathrm{ord}}(q)=\sum_{R\in\mathcal R_{\mathrm{ord}}}
 e^{2\pi iqR}B_R,
 \qquad B_{-R}=B_R^\dagger.
 \label{eq:ordinary-model}
\end{equation}
For an integer cutoff $R_c$, the symmetric truncation is
\begin{equation}
 B^{(R_c)}(q)=\sum_{\abs{R}\le R_c}e^{2\pi iqR}B_R,
 \label{eq:truncated-model}
\end{equation}
with ordinary omitted power
\begin{equation}
 P_{>R_c}^{\mathrm{ord}}[B]
 =\sum_{\abs{R}>R_c}\frac12\norm{B_R}_F^2.
 \label{eq:ordinary-omitted-power}
\end{equation}
The total omitted-power diagnostic adds the separately retained Nyquist conversion power,
\begin{equation}
 P_{>R_c}^{\mathrm{total}}[B]
 =P_{>R_c}^{\mathrm{ord}}[B]+P_{\mathrm{Nyq}}[B].
 \label{eq:total-omitted-power}
\end{equation}
This addition is a bookkeeping identity for the finite conversion; it is not an
ordinary-lattice hopping coefficient at $R=-N/2$.

Truncation is evaluated on a shifted mesh of size $M=1025$,
\begin{equation}
 \widetilde q_\ell=-\frac12+\frac{\ell+\omega}{M},
 \qquad
 \omega=\frac{1}{N+1},
 \qquad \ell=0,\ldots,M-1,
 \label{eq:shifted-evaluation-mesh}
\end{equation}
which is checked to be disjoint from the construction mesh.  At those points the study
retains four norms for each relevant difference: maximum normalized Frobenius error,
root-mean-square normalized Frobenius error, maximum sorted-band error, and
root-mean-square sorted-band error.  For matrix fields $C_\ell$ and $D_\ell$,
\begin{align}
 E_{H,\max}(C,D)&=\max_\ell\frac{\norm{C_\ell-D_\ell}_F}{\sqrt2},\\
 E_{H,\mathrm{rms}}(C,D)&=
 \sqrt{\frac{1}{2M}\sum_\ell\norm{C_\ell-D_\ell}_F^2},\\
 E_{\lambda,\max}(C,D)&=\max_{\ell,n}
 \abs{\lambda_n(C_\ell)-\lambda_n(D_\ell)},\\
 E_{\lambda,\mathrm{rms}}(C,D)&=
 \sqrt{\frac{1}{M}\sum_\ell\sum_{n=1}^2
 [\lambda_n(C_\ell)-\lambda_n(D_\ell)]^2}.
 \label{eq:truncation-errors}
\end{align}
The implementation checks the finite-dimensional inequality that the maximum sorted-band
error does not exceed the maximum Frobenius error beyond a binary64 roundoff allowance.

Three matrix comparisons remain distinct.  The interpolation channel compares the full
finite comparison Fourier interpolant with a direct gauge action on the interpolated
reference.  The truncation channel compares truncated and complete comparison
interpolants.  The total channel compares the truncated comparison with the direct gauge
reference.  Their errors may be numerically similar in a well-resolved case, but they
are not the same defined quantity.

\subsection{Canonical aggregate and lifecycle}

The calculation retains the parent family, momentum refinement, adaptive amplitude
outcome, and truncation outcome in one immutable typed aggregate.  Channel variants are
closed: no type registry, null padding, or untyped dictionary substitutes for a missing
scientific outcome.  Canonical serialization follows
\begin{equation}
 \texttt{\_mapping}\longrightarrow\texttt{to\_dict}
 \longrightarrow\texttt{to\_json},
 \label{eq:serialization-forward}
\end{equation}
with the inverse path
\begin{equation}
 \texttt{\_from\_mapping}\longrightarrow\texttt{from\_dict}
 \longrightarrow\texttt{from\_json}.
 \label{eq:serialization-inverse}
\end{equation}
The calculation script checks byte-for-byte canonical round trip before atomic
publication.  Retention owns and authenticates those bytes and the canonical
visualization projection.  The strictly downstream renderer admits stored verification
and does not rerun scientific verification.
